\begin{frame}
  \frametitle{\bfseries\sffamily Redes Neuronales Artificiales}

  Matemáticamente hablando, las \alert{redes neuronales artificiales}
  (RNA) son funciones compuestas por bloques de construcción muy
  simples con una gran cantidad de parámetros libres.
  Estos parámetros pueden especificarse para que la red neuronal
  artificial realice una tarea matemática específica, lo que
  significa simplemente que se aproxima a una función que determina
  el resultado de la tarea.
  Por lo tanto, el \alert{aprendizaje automático}, y en particular el
  trabajo con redes neuronales artificiales, puede considerarse un
  tipo de aproximación.
  Sin embargo, a diferencia de la interpolación polinómica o
  trigonométrica, el \alert{aprendizaje automático} con
  \alert{redes neuronales artificiales} utiliza un tipo de función
  completamente diferente.
  En este caso, no se elige un espacio vectorial estructurado, como el
  espacio de polinomios o polinomios trigonométricos, para aproximar
  una función, sino que se ensamblan funciones a partir de componentes
  elementales: las neuronas artificiales.

  \pause

  \

  Inicialmente, la investigación matemática sobre redes neuronales
  artificiales se inspiró en sus contrapartes naturales, es decir, en
  el funcionamiento de las neuronas cerebrales.
  Hoy en día, constituyen una herramienta enormemente poderosa que
  puede utilizarse de forma muy eficiente en los ordenadores modernos.
  Muchos problemas antiguos se han resuelto finalmente con éxito
  gracias a las redes neuronales artificiales.
  Cabe destacar los avances en el procesamiento del lenguaje, como la
  traducción automática, el reconocimiento de voz y los modelos de
  lenguaje generativo como ChatGPT, así como los importantes logros
  en el procesamiento de imágenes y vídeos.

  \pause

  \

  Matemáticamente, podemos considerar las redes neuronales
  artificiales como un método para aproximar funciones.
  Sin embargo, su análisis es más complejo y se carece de muchas
  herramientas: la interpolación de Lagrange se basaba
  principalmente en la expansión de Taylor y la representación del
  polinomio de interpolación con polinomios base.
  Sin embargo, las redes neuronales a menudo no son diferenciables
  ni representan un espacio vectorial con una base.

  \pause

  \

  La mejor aproximación gaussiana es un problema de minimización con
  un funcional cuadrático, que puede resolverse con precisión
  mediante el producto escalar.
  Las redes neuronales también se determinan mediante problemas de
  minimización, pero estas tareas de minimización tienen una
  estructura no lineal compleja, no son cuadráticas y, por lo
  general, no son unívocamente resolubles.
  Por lo tanto, carecen de propiedades estructurales importantes
  para el análisis.
\end{frame}